% Part: incompleteness
% Chapter: representability-in-q
% Section: prim-rec

\documentclass[../../../include/open-logic-section]{subfiles}

\begin{document}

\olfileid[id]{inc}{req}{pri}
\olsection{Menyimulasikan Rekursi Primitif}

Sekarang kita dapat menunjukkan bahwa definisi dengan rekursi primitif dapat
``disimulasikan'' oleh pencarian tak berbatas pada fungsi reguler dengan
menggunakan fungsi beta. Andaikan kita mempunyai $f(\vec x)$ dan
$g(\vec x, y, z)$. Maka fungsi~$h(\vec x,y)$ yang didefinisikan dari $f$
dan~$g$ dengan rekursi primitif ialah
\begin{align*}
h(\vec x, 0) & =  f(\vec x) \\
h(\vec x, y+1) & =  g(\vec x, y, h(\vec x, y)).
\end{align*}
Kita perlu menunjukkan bahwa $h$ dapat didefinisikan dari $f$ dan~$g$ hanya
dengan menggunakan komposisi dan pencarian tak berbatas pada fungsi reguler,
dengan fungsi-fungsi dasar dan fungsi-fungsi yang didefinisikan darinya
melalui komposisi dan pencarian tak berbatas pada fungsi reguler (seperti
$\beta$).

\begin{lem}
\ollabel{lem:prim-rec}
Jika $h$ dapat didefinisikan dari $f$ dan $g$ dengan menggunakan rekursi
primitif, fungsi itu dapat didefinisikan dari $f$, $g$, fungsi-fungsi $\Zero$,
$\Succ$, $\Proj{n}{i}$, $\Add$, $\Mult$, $\Char{=}$, dengan menggunakan
komposisi dan pencarian tak berbatas pada fungsi reguler.
\end{lem}

\begin{proof}
Pertama, definisikan fungsi bantu $\hat h(\vec x, y)$ yang mengembalikan
bilangan~$d$ terkecil sedemikian sehingga $d$ mengodekan suatu barisan yang
memenuhi
\begin{enumerate}
\item $(d)_0 = f(\vec x)$, dan
\item untuk setiap $i < y$, $(d)_{i+1} = g(\vec x, i, (d)_i)$,
\end{enumerate}
dengan $(d)_i$ sekarang merupakan singkatan bagi $\beta(d,i)$. Dengan kata
lain, $\hat h$ mengembalikan barisan
$\tuple{h(\vec x, 0), h(\vec x, 1), \dots, h(\vec x, y)}$. Kita dapat
menuliskan $\hat h$ sebagai
\[
\hat h(\vec x, y) = \umin{d}{(\beta(d,0) = f(\vec x) \land \bforall{i <
  y}{\beta(d,i+1) = g(\vec x, i,\beta(d,i))})}.
\]
Perhatikan: rekursi primitif tidak diperlukan di sini, hanya pencarian tak berbatas.
Fungsi yang kita minimalkan bersifat reguler menurut lema fungsi beta
\olref[bet]{lem:beta}.

Namun, sekarang kita mempunyai
\[
h(\vec x, y) = \beta(\hat h(\vec x, y), y),
\]
sehingga $h$ dapat didefinisikan dari fungsi-fungsi dasar hanya dengan
menggunakan komposisi dan pencarian tak berbatas pada fungsi reguler.
\end{proof}

\end{document}
